% Pursuing Stacks (cotes 134-2 to 134-8): archive cote 134-8, site page 7 (PDF page 8), typescript page 561.
% typescript. Draft compared: ps7.tex:297-376.
% AI-transcribed by Claude Opus 5.5 (claude-opus-5-5, adaptive thinking, effort high), Message Batches API, batch ps4 (msgbatch_012Aoc1NuskATPsyBeKSYK7J), 2026-11-12; not yet checked by eye.
\scanpage{134-8/7}
\mspage{561}

make any attempt,  now or later,  to come to an answer.  The last one
anyhow seems closely related to the formalism of closed model structures
on categories $A^{\wedge}$,  and the proper place for dealing with it would seem
to be rather next chapter VII,  where such structures are going to be
\xout{introduced} studied.  What I would like to do however here-and-now,  is
to establish at last the long promised relationship ``integrators are
abelianizators'',  which we have kept t\unsure{ur}ning around and postponing
ever since section 99,  when those integrators were finally introduced
mainly for this purpose (of furnishing us with ``abelianizators'').

\inserted{First of all,} I should be more outspoken than I have been before,  \xout{first of} in
defining the ``\uline{abelianization functor}'' \inserted{(or ``\uline{absolute Whitehead functor}''):}
\[ (1) \qquad \hand{\uline{Wh}} \ : \ (\mathrm{Hot}\hand{)} \overset{\mathrm{def}}{=} \underline{W}^{-1}(\mathrm{Cat}) \longrightarrow \mathrm{Hotab} \overset{\mathrm{def}}{=} \mathrm{D.(Ab)} \quad , \]
without any use of the semisimplicial machinery\deleted{,} which\hand{,} at the beginning
of our reflections\hand{,} ha\over{s}{d} rather obscured the picture\deleted{,} (section 92).  Defi-%
ning such a functor amounts to defining a functor
\[ (2) \qquad \mathrm{LH.} : \ (\mathrm{Cat}\hand{)} \longrightarrow \mathrm{D.(Ab)} \quad , \]
or ``total homology functor'' ,  which should take weak equivalences into
isomorphisms.  For an\xout{\illegible{1}} object A in (Cat),  we define
\[ (3) \qquad \mathrm{LH.}(A) = \mathrm{LH.}(A,\underline{Z}_A) = \mathrm{L}p^{ab}_{A!}\,(\underline{Z}_A) \ , \]
where $\underline{Z}_A$ is the constant abelian presheaf on A with value $\underline{Z}$ ,  and
\[ p_A : A \longrightarrow e \]
is the map to the final object of (Cat).  We have to define the functo-%
rial dependance on A ,  More generally ,  for pairs
\[ (A,F) \ , \ \text{with } A \text{ in (Cat)} \ , \ F \text{ in } A^{\wedge}_{ab} \quad , \]
the expression
\[ \mathrm{LH.}(A,F) \ \text{ in } \ \mathrm{D.(Ab)} \]
is functorial with respect to the pair (A,F),  where a map
\[ (A,F) \longrightarrow (A',F') \]
is defined to be a pair  \xout{fxxAxxxxxA'xx}  (f,u),  where
\[ (4) \qquad f: A \to A' \ , \ u: F \to u^{\hand{*}}(F') \ , \]
the composition of \inserted[typed]{maps} \xout{\illegible{1}} being the obvious  one.  To see that such
a pair  defines a map
\[ (5) \qquad \mathrm{LH.}(f,u): \ \mathrm{LH.}(A,F) \longrightarrow \mathrm{LH.}(A',F') \ , \]

\margin{7}
